% Drop-in statistical reporting for the current full-task results.
% Generated from analyze_task_statistics.py using independent-group tests.

\paragraph{Statistical analysis.}
Task-level success proportions are reported with two-sided 95\% exact
Clopper--Pearson confidence intervals. Because the three system variants were
evaluated in separate trial sets, pairwise differences were assessed using
two-sided Fisher's exact tests. The primary comparison was the recoverable
system against the proposed single-shot system; comparisons with ACT were
treated as secondary.

The recoverable system achieved a TSR of 92.0\% (46/50; 95\% CI:
80.8--97.8\%), compared with 78.0\% for the single-shot system (39/50; 95\%
CI: 64.0--88.5\%). This corresponds to an observed improvement of 14.0
percentage points (Newcombe 95\% CI for the difference:
$-0.2$--28.1 percentage points). The difference did not reach the conventional
0.05 significance level in a two-sided Fisher's exact test
($p=0.091$). We therefore interpret the result as an observed practical
improvement rather than conclusive statistical evidence of superiority.

Relative to ACT (24/50, 48.0\%; 95\% CI: 33.7--62.6\%), the single-shot and
recoverable systems showed differences of 30.0 percentage points
($p=0.0035$) and 44.0 percentage points ($p<0.001$), respectively. These
cross-method comparisons should be interpreted together with the differences
in policy structure and recovery capability described in the experimental
protocol.

% Optional compact table if confidence intervals do not fit in the main task
% table. Keep the original count-based task table and cite this table after it.
\begin{table}[t]
\caption{Task success rates with exact 95\% binomial confidence intervals.}
\label{tab:task_statistics}
\centering
\small
\begin{tabular}{lcc}
\hline
Method & TSR & 95\% CI \\
\hline
ACT & 48.0\% (24/50) & 33.7--62.6\% \\
Ours single-shot & 78.0\% (39/50) & 64.0--88.5\% \\
Ours recoverable & 92.0\% (46/50) & 80.8--97.8\% \\
\hline
\end{tabular}
\end{table}
